% Irodov Part VI §6.7 original problems 291--300.
% User-supplied Arihant PDF pp.274--275 (printed pp.265--266);
% answers PDF p.360 (printed 351). Particle rest energies are MeV.
\elementaryproblem{291}{265}
质子的动量大小分别为
$0.10,\ 1.0,\ 10\,\mathrm{GeV}/c$，
求三种情形的动能。
\exsolution
质子的静止能量$m_pc^2=0.93827\,\mathrm{GeV}$。
由相对论质能关系
$E^2=p^2c^2+m_p^2c^4$、
$T=E-m_pc^2$，得到
\[
 \boxed{T(p)=\sqrt{p^2c^2+m_p^2c^4}-m_pc^2}.
\]
代入三种动量
\[
 \boxed{T_{0.10}=0.00531\,\mathrm{GeV}
 =5.31\,\mathrm{MeV}},\qquad
 \boxed{T_{1.0}=0.433\,\mathrm{GeV}},\qquad
 \boxed{T_{10}=9.11\,\mathrm{GeV}}.
\]
\excheck 本版书末首项印$0.05\,\mathrm{GeV}$，
但动量仅$0.10\,\mathrm{GeV}/c$时
非相对论检验$T\simeq p^2/(2m_p)=5.3$MeV；
书末首项有约十倍数量级差异，需记录。

\elementaryproblem{292}{265}
$\pi$介子的动能超过其静止能量$\eta=1.2$倍，
即$T=\eta m_\pi c^2$。
固有平均寿命$\tau_0=25.5\,\mathrm{ns}$，
求实验室系平均飞行距离。
\exsolution
洛伦兹因子
$\gamma=1+T/(m_\pi c^2)=1+\eta$。
平均寿命延长至$\gamma\tau_0$，
速率$v=\beta c$，
所以
\[
 \overline l=\beta\gamma c\tau_0
 =c\tau_0\sqrt{\gamma^2-1}
 =\boxed{c\tau_0\sqrt{\eta(\eta+2)}}.
\]
代入$\eta=1.2$得
\[
 \boxed{\overline l\simeq
 (299792458)(25.5\times10^{-9})\sqrt{1.2(3.2)}
 \simeq15.0\,\mathrm m}.
\]
\excheck 既要考虑寿命膨胀$\gamma$，
也要计算运动速度$\beta c$；
不可简单取$c\gamma\tau_0$。

\elementaryproblem{293}{266}
负$\pi$介子动能
$T=100\,\mathrm{MeV}$，
从产生至衰变的平均实验室距离
$l=11\,\mathrm m$。
求其固有平均寿命。
\exsolution
定义$m_\pi c^2\simeq139.57\,\mathrm{MeV}$。
由$pc=\sqrt{T(T+2m_\pi c^2)}$，
$\beta\gamma=pc/(m_\pi c^2)$。
相对论飞行距离$l=\beta\gamma c\tau_0$，
反解
\[
 \boxed{\tau_0=
 \frac{l\,m_\pi c^2}
 {c\sqrt{T(T+2m_\pi c^2)}}}.
\]
代入$T=100$MeV及$l=11$m，
\[
 \boxed{\tau_0\simeq2.63\times10^{-8}\,\mathrm s
 \simeq26\,\mathrm{ns}}.
\]
\excheck 得到的固有寿命与低速$\pi^\pm$
约26ns的寿命一致，说明已正确剥离时间膨胀。

\elementaryproblem{294}{266}
一束单能负$\pi$介子，其动能恰等于静止能量
$T=m_\pi c^2$，
在两处相距$l=20\,\mathrm m$的束流截面测通量。
固有平均寿命$\tau_0=25.5\,\mathrm{ns}$。
求后、前两截面的通量比。
\exsolution
$T=m_\pi c^2$意味着$\gamma=2$，
$\beta\gamma=\sqrt{\gamma^2-1}=\sqrt3$。
平均衰变长度
$L=\beta\gamma c\tau_0=\sqrt3c\tau_0$。
稳定束源条件下通过距离$l$的生存概率
就是通量比：
\[
 \boxed{\frac{J(l)}{J(0)}
 =\exp\!\left(-\frac{l}{\sqrt3 c\tau_0}\right)}.
\]
代入得到
$\boxed{J(l)/J(0)\simeq0.22}$。
\excheck 两个截面之间只有粒子衰变引起的损失；
不计束流几何发散或物质吸收。

\elementaryproblem{295}{266}
静止的正$\pi$介子衰变为
$\mu^+$和中微子，
求$\mu^+$动能及中微子能量。
\exsolution
取$m_\pi c^2=139.57$MeV、
$m_\mu c^2=105.66$MeV，
中微子静质量近似为零。
静止母粒子的二体衰变满足
$\boldsymbol p_\mu=-\boldsymbol p_\nu$。
由$E_\mu+E_\nu=m_\pi c^2$与
$E_\nu=p_\nu c$，
\[
 E_\nu=\frac{m_\pi^2-m_\mu^2}{2m_\pi}c^2,
 \quad
 T_\mu=\frac{(m_\pi-m_\mu)^2}{2m_\pi}c^2.
\]
代入
\[
 \boxed{T_\mu\simeq4.12\,\mathrm{MeV}},\qquad
 \boxed{E_\nu\simeq29.79\,\mathrm{MeV}}.
\]
\excheck 中微子能量与$\mu$子的动能不同，
因为$\mu$子携带显著静止能量。

\elementaryproblem{296}{266}
静止的$\Sigma^-$超子经
$\Sigma^-\to n+\pi^-$衰变，
求末态中子的动能。
\exsolution
取静止能量
$m_\Sigma c^2\simeq1197.45$MeV、
$m_nc^2\simeq939.57$MeV、
$m_\pi c^2\simeq139.57$MeV。
二体衰变的中子总能量
\[
 E_n=\frac{m_\Sigma^2+m_n^2-m_\pi^2}
 {2m_\Sigma}c^2,
\]
因此
\[
 \boxed{T_n=E_n-m_nc^2
 =\frac{(m_\Sigma-m_n)^2-m_\pi^2}
 {2m_\Sigma}c^2
 \simeq19.6\,\mathrm{MeV}}.
\]
\excheck 两粒子等动量而非等动能；
书末近似$19.5$MeV。

\elementaryproblem{297}{266}
一个初始静止的正$\mu$子衰变成
正电子和两个中微子。
求正电子可能具有的最大动能。
\exsolution
三体衰变中，当两中微子近似平行同向、
其组合的不变质量达到最小值零时，
正电子总能量最大。
视“两个中微子”为一个零质量系统，
二体运动学给
\[
 E_{e,\max}=
 \frac{m_\mu^2+m_e^2}{2m_\mu}c^2.
\]
减正电子静能
\[
 \boxed{T_{e,\max}=
 \frac{(m_\mu-m_e)^2}{2m_\mu}c^2
 \simeq52.32\,\mathrm{MeV}}.
\]
\excheck 若忽略电子质量，
近似为$m_\mu c^2/2\simeq52.8$MeV；
书末写约52.5MeV，属于所用质量常数精度差异。

\elementaryproblem{298}{266}
一个初始静止的电中性粒子
衰变为动能$T_p=5.3\,\mathrm{MeV}$的质子
及$\pi^-$。求母粒子的静止质量并指出名称。
\exsolution
静止母粒子衰变后二体动量等大反向，
因此
\[
 pc=\sqrt{T_p(T_p+2m_pc^2)},\qquad
 E_\pi=\sqrt{p^2c^2+m_\pi^2c^4}.
\]
母粒子静能
$Mc^2=m_pc^2+T_p+E_\pi$，
取$m_pc^2=938.27$MeV、
$m_{\pi^-}c^2=139.57$MeV，得
\[
 \boxed{Mc^2\simeq1115.2\,\mathrm{MeV}
 \approx1.115\,\mathrm{GeV}}.
\]
对应粒子为
$\boxed{\Lambda^0\text{ 超子}}$，
其标准静止能量约1115.7MeV。
\excheck 不能把质子动能和$\pi$的静止能直接相加；
$\pi$粒子还携带与质子等模的反冲动量。

\elementaryproblem{299}{266}
负$\pi$介子以动能$T_\pi=50\,\mathrm{MeV}$
飞行时衰变成$\mu$子和中微子。
若中微子在实验室系中与母介子运动方向
成$90^\circ$，求中微子能量。
\exsolution
四动量守恒
$P_\mu=P_\pi-P_\nu$，
结合洛伦兹不变量
$P_\mu^2=m_\mu^2c^2$。
因为中微子与母介子动量垂直，
$P_\pi\cdot P_\nu=E_\pi E_\nu/c^2$，
且$P_\nu^2=0$，所以
\[
 m_\mu^2c^4=m_\pi^2c^4-2E_\pi E_\nu,
 \qquad E_\pi=m_\pi c^2+T_\pi.
\]
解得
\[
 \boxed{E_\nu=
 \frac{(m_\pi^2-m_\mu^2)c^4}
 {2(m_\pi c^2+T_\pi)}
 \simeq21.9\,\mathrm{MeV}\approx22\,\mathrm{MeV}}.
\]
\excheck 实验室垂直出射条件使初态
空间动量点积为零，不能套静止母粒子的29.8MeV。

\elementaryproblem{300}{266}
动能$T_\Sigma=320\,\mathrm{MeV}$的
$\Sigma^+$超子飞行中衰变为
一粒电中性粒子和一粒$\pi^+$。
$\pi^+$以相对母粒子前进方向$90^\circ$出射，
其动能$T_\pi=42\,\mathrm{MeV}$。
求中性产物的静止质量。
\exsolution
由四动量
$P_X=P_\Sigma-P_\pi$，
及两动量在实验室空间中垂直，
得到
\[
 m_X^2c^4=
 m_\Sigma^2c^4+m_\pi^2c^4-2E_\Sigma E_\pi,
\]
\[
 E_\Sigma=m_\Sigma c^2+T_\Sigma,\qquad
 E_\pi=m_\pi c^2+T_\pi.
\]
取$m_{\Sigma^+}c^2\simeq1189.37$MeV、
$m_{\pi^+}c^2\simeq139.57$MeV，
\[
 \boxed{m_Xc^2\simeq
 \sqrt{1189.37^2+139.57^2-
 2(1189.37+320)(139.57+42)}
 \simeq941\,\mathrm{MeV}\approx0.94\,\mathrm{GeV}}.
\]
中性产物为$\boxed{n\text{（中子）}}$。
\excheck 输入值取两位有效数字，
所得约941MeV与中子939.6MeV在近似精度内一致。
